\documentclass{sig-alternate}

\usepackage{times}
\usepackage{url}
\usepackage{todonotes}
\sloppy

\setlength{\parindent}{0.5cm} 

\newcommand{\citep}[1]{\cite{#1}}

% Copyright
\setcopyright{acmcopyright}
%\setcopyright{acmlicensed}
%\setcopyright{rightsretained}
%\setcopyright{usgov}
%\setcopyright{usgovmixed}
%\setcopyright{cagov}
%\setcopyright{cagovmixed}


% DOI
\doi{10.475/123_4}

% ISBN
\isbn{123-4567-24-567/08/06}

\acmPrice{\$15.00}

% You need the command \numberofauthors to handle the 'placement
% and alignment' of the authors beneath the title.
%
% For aesthetic reasons, we recommend 'three authors at a time'
% i.e. three 'name/affiliation blocks' be placed beneath the title.
%
% NOTE: You are NOT restricted in how many 'rows' of
% "name/affiliations" may appear. We just ask that you restrict
% the number of 'columns' to three.
%
% Because of the available 'opening page real-estate'
% we ask you to refrain from putting more than six authors
% (two rows with three columns) beneath the article title.
% More than six makes the first-page appear very cluttered indeed.
%
% Use the \alignauthor commands to handle the names
% and affiliations for an 'aesthetic maximum' of six authors.
% Add names, affiliations, addresses for
% the seventh etc. author(s) as the argument for the
% \additionalauthors command.
% These 'additional authors' will be output/set for you
% without further effort on your part as the last section in
% the body of your article BEFORE References or any Appendices.

\begin{document}

\conferenceinfo{GECCO'16,} {July 20-24, 2016, Denver, Colorado, USA.}
\CopyrightYear{2016}
\crdata{TBA}
\clubpenalty=10000
\widowpenalty = 10000

\title{Visualizing genetic programming ancestries using graph databases}

%\numberofauthors{4}
%\author{
%\alignauthor
%Nicholas Freitag McPhee\\
%	\affaddr{Div. of Science and Math}\\
%	\affaddr{Univ. of Minnesota, Morris}\\
%	\affaddr{Morris, MN USA-56267}\\
%	\email{mcphee@morris.umn.edu}
%\alignauthor
%Maggie M. Casale\\
%	\affaddr{Div. of Science and Math}\\
%	\affaddr{Univ. of Minnesota, Morris}\\
%	\affaddr{Morris, MN USA-56267}\\
%	\email{casal033@morris.umn.edu}
%\and
%\alignauthor
%Mitchell Finzel \\
%	\affaddr{Div. of Science and Math}\\
%	\affaddr{Univ. of Minnesota, Morris}\\
%	\affaddr{Morris, MN USA-56267}\\
%	\email{finze008@morris.umn.edu}
%\and
%\alignauthor
%Thomas Helmuth\\
%	\affaddr{Computer Science Dep't}\\
%	\affaddr{Washington and Lee Univ.}\\
%	\affaddr{Lexington, VA USA-24450}\\
%	\email{helmutht@wlu.edu}
%\alignauthor
%Lee Spector\\
%	\affaddr{Cognitive Science}\\
%	\affaddr{Hampshire College}\\
%	\affaddr{Amherst, MA USA-01002}\\
%	\email{lspector@hampshire.edu}
%}

\maketitle

\begin{abstract}

Previous work has demonstrated the utility of graph databases as a tool for collecting and analyzing ancestry in evolutionary computation runs. That work focused on sections of individual runs, whereas this poster illustrates the application of these ideas on the entirety of large runs (up to one million individuals) and combinations of multiple runs.

Here we use these tools to generate graphs showing \emph{all} the ancestors of successful individuals from a variety of stack-based genetic programming runs on software synthesis problems. These graphs highlight important moments in the evolutionary process. They also allow us to compare the dynamics when using different evolutionary tools, such as different selection mechanisms or representations, as well as comparing the dynamics for successful and unsuccessful runs.

\end{abstract}

\begin{CCSXML}
	<ccs2012>
	<concept>
	<concept_id>10003120.10003145.10003146.10010892</concept_id>
	<concept_desc>Human-centered computing~Graph drawings</concept_desc>
	<concept_significance>500</concept_significance>
	</concept>
	<concept>
	<concept_id>10003120.10003145.10003147.10010364</concept_id>

	<concept_desc>Human-centered computing~Scientific visualization</concept_desc>
	<concept_significance>500</concept_significance>
	</concept>
	<concept>
	<concept_id>10010147.10010178.10010205.10010206</concept_id>
	<concept_desc>Computing methodologies~Heuristic function construction</concept_desc>
	<concept_significance>300</concept_significance>
	</concept>
	<concept>
	<concept_id>10010147.10010257.10010293.10011809.10011813</concept_id>
	<concept_desc>Computing methodologies~Genetic programming</concept_desc>
	<concept_significance>300</concept_significance>
	</concept>
	</ccs2012>
\end{CCSXML}

\ccsdesc[500]{Human-centered computing~Graph drawings}
\ccsdesc[500]{Human-centered computing~Scientific visualization}
\ccsdesc[300]{Computing methodologies~Heuristic function construction}
\ccsdesc[300]{Computing methodologies~Genetic programming}

\printccsdesc

\keywords{visualization; genetic programming; graph database; ancestry}

\section{Introduction}
\label{sec:introduction}

Reporting results of genetic programming (GP) and evolutionary computation is frequently limited to aggregate statistics such as mean best fitness or percentage of successful runs. Unfortunately this fails to convey the complex dynamics of such evolutionary systems and obscures or omits potentially valuable information about \emph{why} the runs behaved as they did. Previous work~\cite{McPhee:2015:GPTP} has demonstrated the utility of graph databases as tools for collecting and analyzing ancestry in GP runs, but was focused on sections of individual runs.

In this poster we illustrate the use of these tools as a means of exploring entire ancestry trees. We use the Titan graph database\footnote{\url{http://thinkaurelius.github.io/titan/}} along with the Gremlin shell and the Tinkerpop query tools\footnote{\url{https://tinkerpop.incubator.apache.org/}} to store the parent-child relationships from genetic programming runs, and to extract the ancestry trees of specified individuals. We then visualize these subgraphs using the Graphviz \texttt{dot} graph layout tool\footnote{\url{http://www.graphviz.org/}}.

\section{Our test environment: PushGP \& Replace-Space-With-Newline}
\label{sec:RSWN}

All the visualizations presented here are on runs using the Clojush
implementation\footnote{https://github.com/lspector/Clojush} of the 
PushGP genetic programming (GP) system~\cite{spector:2002:GPEM, 1068292}, 
but all the ideas here
could easily apply to almost any evolutionary computation system. The use
of graph databases to capture ancestor lineage is a completely general
concept, for example, and could be applied to any system that implements a
notion of descent with modification. Some of the specifics of the visualization
are tied to particular metrics (e.g., the Damerau-Levenshtein distance
between Push program genomes), but these could be replaced with other
metrics as appropriate to the application domain (e.g., manhattan distances
between genetic algorithms bitstrings, or tree edit distances in
tree-based GP).

We have applied the visualizations described here on a 
number of different problems taken from a suite of software synthesis benchmark 
problems~\cite{Helmuth:2015:GECCO}, but in the interest of space we
will focus here on a single test problem: replace-space-with-newline.
In replace-space-with-newline the goal is to evolve a program that takes a string
as input and has two tasks: (a) it should print the input string, with all spaces
replaced by newlines, and (b) it should return an integer that is the number of
characters in the input string that were \emph{not} replaced. The use of
these two distinct types of error values will be used in some of the visualizations
as described in Section~\ref{sec:nodes}. The details of the system and
settings used in these runs are as described in~\cite{Helmuth:2015:GECCO}.

\section{Basic Structure}
\label{sec:basics}

\todo[inline]{Run0 with no coloring of nodes, and plain edges -Maggie}

\section{Edges}
\label{sec:edges}
\todo[inline]{-Maggie}

\section{Nodes}
\label{sec:nodes}
\todo[inline]{Box Sizing -Maggie}
\todo[inline]{RBM Color -Nic}
\todo[inline]{Dual Color -Maggie}

\begin{figure*}
	\begin{center}
		\includegraphics[width=0.9\textwidth]{../Figures/run0_RBM_color_filtered_and_full_30000.pdf}
	\end{center}
	\caption{Unfiltered (left) and filtered (right) versions of a successful lexicase run, both with the same color map.}
	\label{fig:lexRun0FilteredAndFull}
\end{figure*}

\section{Filtering}
\label{sec:filtering}

While the large graphs of full runs can provide an excellent ``big picture'' 
view of the run dynamics, there is \emph{so} much information that it can be
difficult to isolate specific features. We can, however, extract and visualize
subgraphs that focus on specific areas or events in the runs.

One simple example of this sort of filtering that we've already introduced is
looking only at individuals that are in fact ancestors of individuals in the
final generation. Most individuals in evolutionary systems aren't ultimately
ancestors of individuals in the final population; they're evolutionary ``dead ends'',
either because they have no children, or their descendants eventually fail to have
children. All the graphs provided so far have used this filtering to substantially
cut down on the number of visualized nodes. This typically reduces the number of nodes in a
generation (row of the graph) from the original 1,000 to 300 or less, and sometimes down
to a few dozen.

\begin{figure*}[tb]
	\begin{center}
		\includegraphics[width=0.9\textwidth]{../Figures/run1_RBM_color_filtered_and_full_60000.pdf}
	\end{center}
	\caption{Unfiltered (left) and filtered (right) versions of a successful lexicase run, both with the same color map.}
	\label{fig:lexRun1FilteredAndFull}
\end{figure*}


\section{Comparing runs}
\label{sec:comparisons}

% In this section we'll look at multiple successful runs in their filtered
% versions. We'll create a RBM color scheme which incorporates all error vectors
% in the filtered ancestries. This will allow us to see similar individual in 
% other runs. Ex) You may see red at the top of one run, and at the end of another. Which could imply that oine found a solution earlier.
\subsection{Successful Runs}
\begin{figure*}[tb]
	\begin{center}
		\includegraphics[height=0.95\textheight]{../Figures/runs_1_99_6_RBM_color_filtered2_with_space.pdf}
	\end{center}
	\caption{Filtered ancestry graphs of three successful runs (1, 99, and 6), all with the same color map.}
	\label{fig:runs1:99:6:filtered}
\end{figure*}

% In the section we'll look at least one successful and failed run in dual
% coloring. This will allow us to see how individuals are prgressing (or not)
% over time. Having these ancestries in full will allow us to generally comapre
% the generation sizes and get the visual impact of a failed run.

\subsection{Successful vs. Failed Runs}
%\begin{figure*}[tb]
%	\begin{center}
%		\includegraphics[width=0.5\textwidth]{../Figures/runs_92_2_dual_color_full_no_last_gen.pdf}
%	\end{center}
%	\caption{Full ancestry graphs of 2 failed runs (92 and 2), all with the same color map.}
%	\label{fig:runs92:2:filtered}
%\end{figure*}

%\begin{figure*}[tb]
%	\begin{center}
%		\includegraphics[width=0.5\textwidth]{../Figures/runs_92_2_1_99_dual_color_full_small.pdf}
%	\end{center}
%	\caption{Full ancestry graphs of 4 different runs (92, 2, 1, and 99), all with the same color map.}
%	\label{fig:runs92:2:1:99:filtered}
%\end{figure*}

\section{Future work}
\label{sec:futurework}
\todo[inline]{Autoconstruction, Dynamic views would be nice. Using google maps to visualize large graphs. Gephi. Centrality metrics. -Nic}



\todo[inline]{Look at previous VizGEC work. - Nic}


\bibliographystyle{abbrv}
\bibliography{GECCO_2016}

\end{document}